\documentclass[11pt,a4paper]{article}

\usepackage{fontspec}
\setmainfont{DejaVu Serif}[
  BoldFont={DejaVu Serif Bold},
  ItalicFont={DejaVu Serif},
  ItalicFeatures={FakeSlant=0.18},
  BoldItalicFont={DejaVu Serif Bold},
  BoldItalicFeatures={FakeSlant=0.18}
]
\setsansfont{DejaVu Sans}
\usepackage{polyglossia}
\setmainlanguage{english}
\usepackage[a4paper,margin=20.5mm]{geometry}
\usepackage{microtype}
\usepackage{setspace}
\setstretch{1.015}
\usepackage{amsmath,amssymb,amsthm,mathtools,bm}
\usepackage{booktabs,array,tabularx}
\usepackage{enumitem}
\usepackage{xcolor}
\usepackage{tikz}
\usetikzlibrary{arrows.meta,positioning,calc}
\usepackage[unicode,hidelinks]{hyperref}
\usepackage[nameinlink,noabbrev]{cleveref}
\emergencystretch=3em
\allowdisplaybreaks

\hypersetup{
  pdftitle={Conical Intersections as Typed Spectral Singularities under Observation Maps},
  pdfauthor={Stas, Independent Research Program},
  pdfsubject={Spectral discriminants, eigenprojectors, Berry holonomy, nonadiabatic coupling, and finite-resolution observation},
  pdfkeywords={conical intersections, Born-Oppenheimer approximation, spectral projector, Berry phase, nonadiabatic coupling, observation geometry}
}

\definecolor{obsblue}{RGB}{37,83,138}
\definecolor{obsred}{RGB}{150,45,45}
\definecolor{obsgreen}{RGB}{30,115,78}
\definecolor{obsviolet}{RGB}{94,64,145}
\definecolor{lightblue}{RGB}{239,246,253}
\definecolor{lightred}{RGB}{253,242,242}
\definecolor{lightgreen}{RGB}{239,249,244}
\definecolor{lightviolet}{RGB}{246,242,252}

\theoremstyle{definition}
\newtheorem{definition}{Definition}[section]
\newtheorem{model}[definition]{Model}
\theoremstyle{plain}
\newtheorem{theorem}[definition]{Theorem}
\newtheorem{proposition}[definition]{Proposition}
\newtheorem{lemma}[definition]{Lemma}
\newtheorem{corollary}[definition]{Corollary}
\theoremstyle{remark}
\newtheorem{remark}[definition]{Remark}
\newtheorem{warning}[definition]{Boundary}

\newcommand{\R}{\mathbb{R}}
\newcommand{\C}{\mathbb{C}}
\newcommand{\dd}{\mathop{}\!d}
\newcommand{\norm}[1]{\left\lVert #1\right\rVert}
\newcommand{\abs}[1]{\left\lvert #1\right\rvert}
\newcommand{\pair}[2]{\left\langle #1,#2\right\rangle}
\newcommand{\Ocal}{\mathcal O}
\newcommand{\Kcal}{\mathcal K}
\newcommand{\Pcal}{\mathcal P}
\newcommand{\Scal}{\mathcal S}
\newcommand{\Hcal}{\mathcal H}
\newcommand{\Dcal}{\mathcal D}
\newcommand{\Ucal}{\mathcal U}
\newcommand{\I}{\mathrm i}
\newcommand{\Tr}{\operatorname{Tr}}
\newcommand{\rank}{\operatorname{rank}}
\newcommand{\Law}{\operatorname{Law}}
\newcolumntype{L}[1]{>{\raggedright\arraybackslash}p{#1}}
\newcolumntype{Y}{>{\raggedright\arraybackslash}X}

\title{\bfseries Conical Intersections as Typed Spectral Singularities\\
\large Eigenprojectors, holonomy, nonadiabatic geometry, and observation maps}
\author{Stas\\
\small Independent Research Program; ORCID: 0009-0000-2294-705X}
\date{Strict GO Core reference revision v1.1 --- 28 July 2026}

\begin{document}
\maketitle

\begin{center}
\fcolorbox{obsblue}{lightblue}{%
\parbox{0.92\linewidth}{\small
\textbf{Status.} This document supersedes the bilingual working note
v1.0. It is an application-layer reference migrated to GO Core v0.2.
It does not propose a new electronic-structure method, nonadiabatic
dynamics algorithm, conical-intersection theorem, or optical-caustic
equivalence. The internal contribution is a typed separation of
nuclear coordinates, energy-valued branching coordinates,
gauge-invariant projectors, local eigenvector sections, dynamics, and
finite-resolution observation.}}
\end{center}

\begin{abstract}
Let \(q\) denote physical nuclear coordinates and let
\(H_{\rm e}(q)\) be an electronic Hamiltonian family with an isolated
two-state spectral cluster. Near a generic real two-state degeneracy,
the linear branching-plane model is
\[
H_{\rm eff}(q)=E_0(q)I+x(q)\sigma_z+y(q)\sigma_x,\qquad
x=g_\mu\delta q^\mu,\quad y=h_\mu\delta q^\mu,
\]
where \(x,y\) are energies, not nuclear lengths.  The seam is the
regular preimage \(x=y=0\), hence has codimension two when
\(\rank D_q(x,y)=2\).  On the punctured branching plane,
\[
E_\pm=E_0\pm r,\quad
P_\pm=\frac12\left(I\pm\frac{x\sigma_z+y\sigma_x}{r}\right),
\quad r=\sqrt{x^2+y^2}.
\]
The unordered spectrum and the two-state cluster remain meaningful at
the degeneracy, but neither rank-one projector has a unique limit
there. Eigenvectors are only gauge-dependent local sections.  A loop
of winding number \(w\) has real-line holonomy
\((-1)^w\) and Berry phase \(\pi w\) modulo \(2\pi\).  In the energy
coordinates, the gauge-invariant quantum metric is
\[
g^{(E)}=\frac1{4r^4}
\begin{pmatrix}y^2&-xy\\-xy&x^2\end{pmatrix},
\qquad \Tr g^{(E)}=\frac1{4r^2}.
\]
Its pullback by \(D_q(x,y)\) supplies the correctly dimensioned
physical-coordinate coupling. A closing gap invalidates a uniform
adiabatic estimate, but does not determine a transition probability;
that requires a nuclear path, speed, initial state, and dynamical
model. Finally, finite spectral resolution can certify only an
unresolved near-degeneracy unless uncertainty and identifiability
conditions are supplied. The strict object is a spectral-discriminant
singularity. ``Observation caustic'' is retained only as a qualified
program-level analogy.
\end{abstract}

\tableofcontents

\section{Scope and typed setting}

\subsection{Electronic family and spectral cluster}

Let \(U\subset\R^f\) be a nuclear-coordinate chart.  For the strict
finite-dimensional reference, \(H_{\rm e}:U\to\operatorname{Herm}(n)\)
is \(C^2\).  The two electronic levels of interest form a cluster
separated from the remaining spectrum by a positive external gap on a
neighborhood \(V\subset U\).  The Riesz projector onto that cluster is
\begin{equation}\label{eq:cluster-projector}
P_{\rm cl}(q)=\frac{1}{2\pi\I}\oint_\Gamma
(z-H_{\rm e}(q))^{-1}\,\dd z,
\end{equation}
where one fixed contour \(\Gamma\) encloses precisely the two cluster
eigenvalues.  Under these hypotheses \(P_{\rm cl}\) is \(C^2\), even
when the two eigenvalues cross internally.

For an unbounded electronic operator, a common dense domain,
self-adjointness, resolvent regularity, and an isolated spectral
cluster must be stated separately.  The finite matrix reference does
not silently prove those operator-theoretic hypotheses.

\subsection{Physical and energy-valued coordinates}

Choose \(q_\ast\in U\), set \(\delta q=q-q_\ast\), and write the
linear two-state reduction after a fixed local basis choice as
\begin{equation}\label{eq:linear-normal-form}
H_{\rm eff}(q)
=E_0(q)I+x(q)\sigma_z+y(q)\sigma_x+O(\norm{\delta q}^2),
\qquad
x=g_\mu\delta q^\mu,\quad y=h_\mu\delta q^\mu.
\end{equation}
Here \(q^\mu\) is a physical nuclear coordinate, while
\(g_\mu,h_\mu\) are energy-gradient covectors. Thus
\[
[q^\mu]=L,\qquad [g_\mu]=[h_\mu]=EL^{-1},\qquad [x]=[y]=E
\]
in the declared length chart. A mass-weighted or curvilinear chart is
allowed only with its own coordinate dimensions and Jacobian.

\begin{warning}[Coordinate firewall]
The notation \(H=Q_1\sigma_z+Q_2\sigma_x\) is dimensionally valid only
if \(Q_1,Q_2\) are declared energy-valued branching coordinates. It
is not valid when the same symbols denote unscaled nuclear lengths.
\end{warning}

\section{Spectral discriminant and codimension}

The exact benchmark drops \(E_0I\) and higher-order terms:
\begin{equation}\label{eq:benchmark-hamiltonian}
H(x,y)=x\sigma_z+y\sigma_x
=\begin{pmatrix}x&y\\y&-x\end{pmatrix},
\qquad
r=\sqrt{x^2+y^2}.
\end{equation}

\begin{proposition}[Conical spectrum]\label{prop:conical-spectrum}
For \((x,y)\ne(0,0)\),
\begin{equation}\label{eq:eigenvalues-gap}
E_\pm=\pm r,\qquad \Delta:=E_+-E_-=2r.
\end{equation}
The repeated eigenvalue occurs precisely at \(x=y=0\).
\end{proposition}

\begin{proof}
\(\det(H-EI)=E^2-x^2-y^2\), so \(E=\pm r\). Equality of the two
eigenvalues is equivalent to \(r=0\).
\end{proof}

Let \(F=(x,y):U\to\R^2\) and \(S=F^{-1}(0)\).

\begin{proposition}[Generic real seam]\label{prop:codimension}
If \(\rank D_qF(q_\ast)=2\), then \(S\) is a \(C^1\) submanifold of
codimension two near \(q_\ast\). Its tangent space is
\[
T_{q_\ast}S=\ker D_qF(q_\ast).
\]
\end{proposition}

\begin{proof}
This is the regular-value theorem applied to \(F\).
\end{proof}

\begin{warning}[Symmetry class matters]
A general complex Hermitian two-state matrix contains three
independent traceless coefficients; its generic degeneracy has
codimension three. Symmetry-enforced crossings, spin-orbit coupling,
Kramers degeneracy, and higher multiplicity require different
discriminants. Codimension two is asserted only for the declared real
two-state class with a transverse parameter map.
\end{warning}

\section{Projectors, eigenlines, and gauge}

On \(\R^2\setminus\{0\}\), define the gauge-invariant rank-one
projectors
\begin{equation}\label{eq:rank-one-projectors}
P_\pm(x,y)
=\frac12\left(I\pm\frac{x\sigma_z+y\sigma_x}{r}\right).
\end{equation}

\begin{proposition}[No rank-one extension through the seam]
\label{prop:no-projector-extension}
Neither \(P_+\) nor \(P_-\) has a continuous extension to
\((x,y)=(0,0)\).  The cluster projector
\(P_{\rm cl}=P_++P_-=I_2\) does extend smoothly.
\end{proposition}

\begin{proof}
Along \(y=0\), the limits of \(P_-\) as \(x\to0^+\) and \(x\to0^-\)
are respectively
\(\operatorname{diag}(0,1)\) and
\(\operatorname{diag}(1,0)\). They differ. The sum is identically
\(I_2\).
\end{proof}

Write \(x=r\cos\phi\), \(y=r\sin\phi\). One real local lift is
\begin{equation}\label{eq:real-eigenvectors}
u_-(\phi)=
\begin{pmatrix}-\sin(\phi/2)\\ \cos(\phi/2)\end{pmatrix},
\qquad
u_+(\phi)=
\begin{pmatrix}\cos(\phi/2)\\ \sin(\phi/2)\end{pmatrix}.
\end{equation}
It satisfies \(u_\pm(\phi+2\pi)=-u_\pm(\phi)\).  A local eigenvector
is not the output of a global gauge-invariant observation map.
Under
\[
u_a(q)\longmapsto e^{\I\chi_a(q)}u_a(q),
\]
the projector \(P_a=\lvert u_a\rangle\langle u_a\rvert\) is unchanged.

\begin{center}
\begin{tikzpicture}[
  node distance=5mm and 9mm,
  box/.style={draw=obsblue,rounded corners,fill=lightblue,
    align=center,inner sep=4pt,font=\scriptsize},
  local/.style={draw=obsviolet,rounded corners,fill=lightviolet,
    align=center,inner sep=4pt,font=\scriptsize},
  arr/.style={-{Latex[length=2mm]},thick,draw=obsblue},
  darr/.style={-{Latex[length=2mm]},thick,dashed,draw=obsviolet}
]
\node[box] (family) {\(H_{\rm e}(q)\)};
\node[box,right=of family] (cluster) {\(P_{\rm cl}(q)\)\\isolated cluster};
\node[box,right=of cluster] (split) {\((E_\pm,P_\pm)\)\\only off \(S\)};
\node[local,below=of split] (section) {\(u_\pm^{(\alpha)}\)\\local gauge section};
\draw[arr] (family) -- (cluster);
\draw[arr] (cluster) -- node[above,font=\scriptsize] {\(\Delta>0\)} (split);
\draw[darr] (split) -- node[right,font=\scriptsize] {gauge choice} (section);
\end{tikzpicture}
\end{center}

\section{Berry holonomy without a global eigenvector}

Let \(C:[0,1]\to\R^2\setminus\{0\}\) be a closed loop and let
\(w(C,0)\in\mathbb Z\) be its winding number around the origin.

\begin{proposition}[Real-line holonomy]\label{prop:berry-holonomy}
For the model \cref{eq:benchmark-hamiltonian}, parallel transport of
either real eigenline around \(C\) has holonomy
\begin{equation}\label{eq:berry-phase}
\operatorname{Hol}(C)=(-1)^{w(C,0)},\qquad
\gamma(C)=\pi w(C,0)\pmod{2\pi}.
\end{equation}
\end{proposition}

\begin{proof}
The polar angle changes by \(2\pi w\). Equation
\cref{eq:real-eigenvectors} therefore changes the real lift by
\((-1)^w\). The ray returns to itself, while the lifted line has the
stated holonomy.
\end{proof}

A gauge-invariant discrete diagnostic for a sufficiently fine loop
\(q_0,\ldots,q_{N-1}\), with nonzero neighboring overlaps, is
\begin{equation}\label{eq:discrete-wilson-loop}
W_N=\prod_{j=0}^{N-1}
\frac{\langle u_j,u_{j+1}\rangle}
{\abs{\langle u_j,u_{j+1}\rangle}},
\qquad u_N=u_0,\qquad
\gamma_N=\arg W_N.
\end{equation}
Local phase factors telescope in the closed product.  In a real gauge,
\(W_N=-1\) for odd winding and \(+1\) for even winding.

\begin{warning}[Bundle and curvature boundary]
Over a punctured plane, a complex eigenline can be represented by a
single-valued complex section, but the Berry connection may have
nontrivial holonomy. In the real model, the real line over a loop of
odd winding is the nontrivial real line bundle. The smooth Berry
curvature vanishes away from the origin in this two-dimensional real
model; the phase must not be attributed to an ordinary nonzero local
curvature density on the punctured plane.
\end{warning}

\section{Derivative coupling and quantum metric}

For local orthonormal sections, define
\[
d_{ab}^{(q)}=\langle u_a,\nabla_q u_b\rangle .
\]
For \(a\ne b\), a gauge change multiplies \(d_{ab}\) by a phase;
therefore \(d_{ab}\) itself is gauge covariant, while its norm and the
quantum metric are gauge invariant.

In the energy coordinates \(\xi=(x,y)\), the gauge
\cref{eq:real-eigenvectors} gives
\begin{equation}\label{eq:energy-nac}
d_{+-}^{(E)}
=-\frac12\nabla_\xi\phi
=\frac1{2r^2}\begin{pmatrix}y\\-x\end{pmatrix},
\qquad
\norm{d_{+-}^{(E)}}=\frac1{2r}.
\end{equation}
Its dimension is \(E^{-1}\), not \(L^{-1}\).

The rank-one quantum metric is
\begin{equation}\label{eq:quantum-metric}
g^{(E)}_{\mu\nu}
=\operatorname{Re}\langle\partial_\mu u_-,
(I-P_-)\partial_\nu u_-\rangle
=\frac12\Tr(\partial_\mu P_-\partial_\nu P_-)
=\frac1{4r^4}
\begin{pmatrix}
y^2&-xy\\
-xy&x^2
\end{pmatrix}_{\mu\nu}.
\end{equation}
Hence
\begin{equation}\label{eq:metric-trace}
\Tr g^{(E)}=\frac1{4r^2}
=\norm{d_{+-}^{(E)}}^2.
\end{equation}

Let \(J=D_q(x,y)\), with entries of dimension \(EL^{-1}\) in the
declared length chart. The physical-coordinate quantities are the
pullbacks
\begin{equation}\label{eq:physical-pullback}
d_{+-}^{(q)}=J^\mathsf T d_{+-}^{(E)},\qquad
g^{(q)}=J^\mathsf T g^{(E)}J,
\end{equation}
so \([d_{+-}^{(q)}]=L^{-1}\) and \([g^{(q)}]=L^{-2}\).
The scalar \(1/(2r)\) is therefore not a coordinate-independent
physical coupling magnitude; it belongs to the Euclidean metric of
the chosen energy-coordinate plane.

\section{Born--Oppenheimer reduction and dynamics firewall}

For a local adiabatic expansion
\[
\Psi(q,r_{\rm e})=\sum_a\chi_a(q)u_a(q;r_{\rm e}),
\]
the nuclear kinetic operator contains
\begin{equation}\label{eq:kinetic-couplings}
(T_N)_{ab}
=-\sum_A\frac{\hbar^2}{2M_A}
\left[
\delta_{ab}\nabla_A^2
+2d_{ab}^{(A)}\cdot\nabla_A
+\tau_{ab}^{(A)}
\right],
\qquad
\tau_{ab}^{(A)}=\langle u_a,\nabla_A^2u_b\rangle .
\end{equation}
Under a coordinate-dependent unitary change within the retained
cluster, covariance belongs to the complete matrix operator.
Discarding the off-diagonal terms is an approximation whose uniform
validity fails near an internal gap closure.

For a prescribed path \(q(t)\), a local adiabatic diagnostic is
\begin{equation}\label{eq:adiabatic-parameter}
\eta_{ab}(t)
=\frac{\hbar\abs{\langle u_a,\dot H_{\rm e}u_b\rangle}}
{\abs{E_b-E_a}^2}.
\end{equation}
The condition \(\eta_{ab}\ll1\) is a model-dependent sufficiency
diagnostic, not a universal error theorem. At a gap closure, the
uniform-gap hypothesis of standard adiabatic estimates fails.

\begin{model}[Landau--Zener control]\label{model:lz}
For
\[
H_{\rm LZ}(t)=vt\,\sigma_z+b\,\sigma_x,\qquad
v>0,\quad b>0,
\]
with evolution from \(t=-\infty\) to \(+\infty\), the exact diabatic
survival probability is
\begin{equation}\label{eq:lz}
P_{\rm D}=\exp\left(-\frac{\pi b^2}{\hbar v}\right).
\end{equation}
\end{model}

This control has minimum gap \(2b\), but changing the sweep rate \(v\)
changes \(P_{\rm D}\). Thus the static gap, Berry phase, or derivative
coupling alone does not determine a molecular branching ratio.
Wavepacket shape, masses, path, velocity, decoherence, additional
states, and the measurement channel are separate inputs.

\section{Observation chain and finite-resolution identifiability}

The strict factorization is
\begin{equation}\label{eq:observation-chain}
(H_{\rm mol},\Psi)
\xrightarrow{\Scal_{\rm cl}}
(H_{\rm cl}(q),\chi(q))
\xrightarrow{\Pcal_{\rm ad}}
(E_\pm(q),P_\pm(q),\chi_\pm(q))
\xrightarrow{\Ocal_\lambda}
Z_\lambda
\xrightarrow{\Kcal_\lambda}
\Law(Y_\lambda\mid Z_\lambda)
\xrightarrow{\widehat\Theta_\lambda}
\widehat\Theta .
\end{equation}
The rank-one adiabatic splitting \(\Pcal_{\rm ad}\) is defined only
off the seam. A local eigenvector section is an additional gauge
choice, not an observation.  A spectroscopic or product channel
\(\Ocal_\lambda\) does not generally reveal the full projector field
or Berry holonomy.

\begin{definition}[Spectral-discriminant singularity]
Let \(\Dcal\subset\operatorname{Herm}(n)\) be the repeated-eigenvalue
discriminant. A conical intersection of the declared real two-state
class is a transverse point or seam in
\[
S=H_{\rm eff}^{-1}(\Dcal)
\]
for which the ordered rank-one eigenprojector map fails to extend
through \(S\) and the leading eigenvalue splitting is conical.
\end{definition}

This is not the same as requiring \(\rank DH_{\rm eff}\) to drop.
The singular map is the eigen-decomposition at the discriminant, not
necessarily the parameter-to-Hamiltonian map.  For a classical
optical caustic, the reusable criterion is rank loss of a ray
projection; a determinant criterion is valid only in a square local
chart.

\begin{warning}[Qualified use of ``observation caustic'']
The term may be used as a GO analogy for singularity of a reduced
description. It is not a theorem that conical intersections and
optical caustics are equivalent singularities, and no Arnold
classification of Lagrangian-map caustics is imported here.
\end{warning}

Suppose measured energies are
\[
Y_\pm=E_\pm+\varepsilon_\pm,\qquad
\mathbb E[\varepsilon_\pm]=0.
\]
Then the measured gap
\(\widehat\Delta=Y_+-Y_-\) has variance
\begin{equation}\label{eq:gap-uncertainty}
\sigma_\Delta^2
=\sigma_+^2+\sigma_-^2
-2\operatorname{Cov}(\varepsilon_+,\varepsilon_-).
\end{equation}
With instrumental resolution \(\epsilon_E\ge0\), a rule such as
\begin{equation}\label{eq:unresolved-rule}
\abs{\widehat\Delta}
\le \kappa\sigma_\Delta+\epsilon_E
\end{equation}
classifies an \emph{unresolved near-degeneracy}. It does not prove an
exact degeneracy, conical topology, odd Berry holonomy, or a
nonadiabatic transition. Those conclusions require additional
sampling, model class, path, and identifiability assumptions.

\section{Gap-opening control and reproducible benchmarks}

The complex Hermitian control
\begin{equation}\label{eq:gapped-control}
H_\delta(x,y)=x\sigma_z+y\sigma_x+\delta\sigma_y,
\qquad \delta\ne0,
\end{equation}
has
\begin{equation}\label{eq:gapped-spectrum}
E_\pm^\delta=\pm\sqrt{x^2+y^2+\delta^2},\qquad
\Delta_\delta=2\sqrt{r^2+\delta^2}\ge2\abs{\delta}.
\end{equation}
It opens a gap on the \(x,y\) plane by leaving the real-symmetric
class. It is therefore a symmetry-breaking control, not the same
real conical-intersection model with a harmless numerical
regularizer.

For the positively oriented circle \(x=R\cos\phi\),
\(y=R\sin\phi\), and the orientation convention used in the release,
the lower-state phase modulo \(2\pi\) has magnitude
\begin{equation}\label{eq:gapped-berry}
\gamma_\delta(R)
=\pi\left(1-\frac{\delta}{\sqrt{R^2+\delta^2}}\right).
\end{equation}
It tends to \(\pi\) as \(\delta\to0^+\) and to zero as
\(\delta/R\to+\infty\); unlike the real sign holonomy, it is not
quantized for \(\delta\ne0\).

\begin{table}[ht]
\centering
\small
\begin{tabularx}{\linewidth}{L{0.31\linewidth} L{0.22\linewidth} Y}
\toprule
Benchmark & Exact value & Purpose \\
\midrule
\(H(3,4)\), energy units & \(E_\pm=\pm5,\ \Delta=10\) &
Conical spectrum and dimensions \\
\(r=0.25\), energy units & \(\norm{d_{+-}^{(E)}}=2\) &
Energy-coordinate coupling scaling \\
Unit circle, winding \(1\) & \(W_N=-1,\ \gamma=\pi\) &
Gauge-invariant discrete holonomy \\
Unit circle, winding \(2\) & \(W_N=+1,\ \gamma=0\) mod \(2\pi\) &
Parity of real-line holonomy \\
\(R=1,\delta=1\) & \(\gamma_\delta=\pi(1-1/\sqrt2)\) &
Nonquantized gapped control \\
\(\hbar v=b^2\) & \(P_{\rm D}=e^{-\pi}\) &
Dimensionless dynamical exponent \\
\bottomrule
\end{tabularx}
\caption{Reference values implemented by the P8 regression suite.}
\end{table}

The executable release additionally checks projector idempotence and
orthogonality, absence of a rank-one limit, cluster smoothness,
gauge invariance of the overlap product, quantum-metric identities,
coordinate pullback dimensions, the Landau--Zener formula, covariance
in \cref{eq:gap-uncertainty}, and zero-gap guards.

\section{Protocol requirements and claim register}

A complete conical-intersection observation record states at least:
\begin{enumerate}[leftmargin=*,itemsep=2pt]
\item nuclear-coordinate chart, coordinate dimensions, reference
      geometry, and frame;
\item electronic Hamiltonian class, domain assumptions, basis,
      symmetry class, and retained spectral cluster;
\item external cluster gap, active-state multiplicity, and
      degeneracy/discriminant convention;
\item physical-to-branching Jacobian \(D_q(x,y)\) and its units;
\item whether outputs are unordered eigenvalues, ordered energies,
      projectors, local eigenvectors, or a diabatic matrix;
\item gauge convention, loop orientation, winding rule, and overlap
      threshold for a discrete Berry phase;
\item nuclear path, speed, masses, initial state, and dynamical
      approximation for every transition claim;
\item spectroscopic or product observation map, resolution,
      stochastic kernel, covariance, and detection threshold;
\item estimator, uncertainty statement, sampling region, and
      identifiability class;
\item explicit boundary between exact normal-form results,
      asymptotic diagnostics, numerical evidence, and chemical
      interpretation.
\end{enumerate}

\begin{table}[ht]
\centering
\small
\begin{tabularx}{\linewidth}{L{0.34\linewidth} L{0.19\linewidth} Y}
\toprule
Claim & Status & Scope \\
\midrule
Double-cone spectrum &
Proposition &
Exact real two-state normal form \\
Codimension-two seam &
Proposition &
\(\rank D_q(x,y)=2\), real two-state class \\
No rank-one projector extension &
Proposition &
Punctured branching plane approaching the origin \\
Berry sign holonomy &
Proposition &
Closed loop with declared winding, real eigenline \\
\(\Tr g^{(E)}=1/(4r^2)\) &
Proposition &
Euclidean energy-coordinate metric \\
Gap closure fixes transition probability &
Rejected &
Dynamics and initial data are additionally required \\
Finite unresolved gap proves a CI &
Rejected &
Resolution alone does not identify topology \\
CI equals an optical caustic &
Rejected &
Only a qualified observation-reduction analogy \\
New quantum-chemical method or theorem &
Rejected &
Not claimed \\
\bottomrule
\end{tabularx}
\end{table}

\begin{center}
\fcolorbox{obsgreen}{lightgreen}{%
\parbox{0.88\linewidth}{\small
The strict program-level conclusion is that a conical intersection is
a transverse encounter with a spectral discriminant for which the
individual adiabatic rank-one splitting becomes singular while the
retained two-state cluster can remain regular. Gauge, dynamics, and
measurement are separate typed layers.}}
\end{center}

\section*{References}
\addcontentsline{toc}{section}{References}
\begingroup
\raggedright
\begin{enumerate}[label={[\arabic*]},leftmargin=*,itemsep=3pt]
\item M. Born and R. Oppenheimer, ``Zur Quantentheorie der Molekeln,''
\emph{Annalen der Physik} 389, 457--484 (1927),
\href{https://doi.org/10.1002/andp.19273892002}{doi:10.1002/andp.19273892002}.
\item H. C. Longuet-Higgins, U. Oepik, M. H. L. Pryce, and
R. A. Sack, ``Studies of the Jahn--Teller effect. II. The dynamical
problem,'' \emph{Proceedings of the Royal Society A} 244, 1--16
(1958),
\href{https://doi.org/10.1098/rspa.1958.0022}{doi:10.1098/rspa.1958.0022}.
\item G. Herzberg and H. C. Longuet-Higgins, ``Intersection of
potential energy surfaces in polyatomic molecules,''
\emph{Discussions of the Faraday Society} 35, 77--82 (1963),
\href{https://doi.org/10.1039/DF9633500077}{doi:10.1039/DF9633500077}.
\item C. A. Mead and D. G. Truhlar, ``On the determination of
Born--Oppenheimer nuclear motion wave functions including
complications due to conical intersections and identical nuclei,''
\emph{Journal of Chemical Physics} 70, 2284--2296 (1979),
\href{https://doi.org/10.1063/1.437734}{doi:10.1063/1.437734}.
\item M. V. Berry, ``Quantal phase factors accompanying adiabatic
changes,'' \emph{Proceedings of the Royal Society A} 392, 45--57
(1984),
\href{https://doi.org/10.1098/rspa.1984.0023}{doi:10.1098/rspa.1984.0023}.
\item C. Zener, ``Non-adiabatic crossing of energy levels,''
\emph{Proceedings of the Royal Society A} 137, 696--702 (1932),
\href{https://doi.org/10.1098/rspa.1932.0165}{doi:10.1098/rspa.1932.0165}.
\item T. Kato, \emph{Perturbation Theory for Linear Operators},
Springer, corrected printing of the second edition, 1995.
\item D. R. Yarkony, ``Diabolical conical intersections,''
\emph{Reviews of Modern Physics} 68, 985--1013 (1996),
\href{https://doi.org/10.1103/RevModPhys.68.985}{doi:10.1103/RevModPhys.68.985}.
\end{enumerate}
\endgroup

\end{document}
